\section{Related Work}

Existing research on cybersecurity for distributed computing infrastructures has primarily focused on isolated domains such as cloud security, edge computing protection, and IoT security frameworks. However, limited work addresses the holistic security challenges of the integrated Cloud--Edge--IoT continuum.

Recent initiatives, including the European Union Cybersecurity Strategy, ENISA reports, and Horizon Europe projects, emphasize the need for sovereignty-driven security architectures. Prior studies on Zero Trust Architectures (ZTA), Post-Quantum Cryptography (PQC), and confidential computing provide important building blocks but lack a unified prioritization and transition methodology.

This work differentiates itself by proposing a structured, end-to-end methodological pipeline that connects architectural capabilities with policy-driven prioritization and phased implementation strategies.

% =========================================================================
% SECTION 3: RELATED WORK AND METHODOLOGICAL POSITIONING
% For IEEE Access — academic register, full citations
% Version: v1.0
% =========================================================================

\section{Related Work and Methodological Positioning}
\label{sec:related_work}

This section surveys the three research traditions whose synthesis
motivates the methodology presented in Section~\ref{sec:methodology}:
technology roadmapping (TRM), capability-based planning (CBP), and
cybersecurity capability maturity assessment. It then reviews the
methodological approaches of all major European cybersecurity roadmap
deliverables and establishes the specific gap that the present work
addresses.

% -----------------------------------------------------------------
% 3.1 TECHNOLOGY ROADMAPPING FOUNDATIONS
% -----------------------------------------------------------------
\subsection{Technology Roadmapping Foundations}
\label{sec:rw_trm}

Technology roadmapping provides the procedural backbone for structured
strategic planning under technological uncertainty. The Cambridge school,
anchored by Phaal, Farrukh, and Probert~\cite{phaal2004}, formalises
roadmapping as a series of facilitated workshops progressing through
\textit{Market $\rightarrow$ Product $\rightarrow$ Technology $\rightarrow$
Charting}, with each workshop gated by the outputs of its predecessor.
The accompanying T-Plan workbook~\cite{phaal2001_tplan} operationalises
this into an industry-applicable fast-start process. The parallel Sandia
tradition, codified by Garcia and Bray~\cite{garcia1997}, structures
roadmapping as a three-phase process (\textit{Preliminary activity
$\rightarrow$ Development $\rightarrow$ Follow-up}) with seven
development steps progressing from product focus through critical
requirements, technology areas, drivers, alternatives, and timelines
to recommendations. Kostoff and Schaller~\cite{kostoff2001} provide
the first formal taxonomy of roadmapping approaches in the engineering
management literature, distinguishing expert-based from bibliometric
methods and defining the node-link structure inherited by subsequent
work.

The canonical TRM vocabulary organises content into three layers:
\textit{drivers} (why), \textit{products/services/capabilities} (what),
and \textit{technologies/resources} (how), plotted against time on a
generic chart~\cite{phaal2004,kerr2022}. Critically, ``capability'' is
an optional middle-layer label in Phaal's formulation, not an organising
primitive with a formal definition. Terms such as ``focus areas'' and
``priorities'' are derivative: priorities are rankings applied to drivers
or gaps, and focus areas are thematic clusters used loosely in
government and cybersecurity contexts but absent from the Cambridge and
Sandia lexicon.

Two observations are significant for the present work. First, no
canonical TRM framework separates ``capability identification'' from
``capability analysis'' as distinct named stages. Garcia and
Bray~\cite{garcia1997} fold both into their Development phase; T-Plan
separates identification (workshops~1--3) from integration (workshop~4)
but labels both as ``development.'' The explicit separation of these
functions into named pipeline stages is a methodological design choice
justified in Section~\ref{sec:methodology}. Second, quantitative
maturity within TRM relies predominantly on Technology Readiness Levels
(TRL)~\cite{mankins1995,iso16290}, with extensions including
Manufacturing Readiness Levels (MRL) and System Readiness Levels
(SRL)~\cite{sauser2008}. Cambridge TRM is predominantly qualitative;
TRL-annotated roadmaps are more common in defence and aerospace
applications. In cybersecurity specifically, quantitative maturity
typically employs Maturity Indicator Levels (MIL~0--3) in the DOE
Cybersecurity Capability Maturity Model (C2M2)~\cite{doe_c2m2} or
Tiers~1--4 in the NIST Cybersecurity Framework~\cite{nist_csf2},
rather than TRL.

Stage-transition logic in TRM draws on Cooper's Stage-Gate
systems~\cite{cooper1990}, which define gates as quality-control
checkpoints comprising deliverables, must-meet and should-meet criteria,
and Go/Kill/Hold/Recycle outputs. Gerdsri, Vatananan, and
Dansamasatid~\cite{gerdsri2009} provide the most operational transition
logic in dynamic TRM, gating three stages (Initiation $\rightarrow$
Development $\rightarrow$ Integration) by explicit criteria: stakeholder
acceptance, completion of information collection, and incorporation into
ongoing planning. Carvalho, Fleury, and Lopes~\cite{carvalho2013}
identify recurrent quality axes for roadmap assessment---completeness,
internal consistency, traceability, feasibility, and stakeholder
legitimacy---that inform the gate criteria adopted in the present
methodology.

% -----------------------------------------------------------------
% 3.2 CAPABILITY-BASED PLANNING
% -----------------------------------------------------------------
\subsection{Capability-Based Planning in Security Contexts}
\label{sec:rw_cbp}

Capability-based planning (CBP) originates in defence doctrine. The
canonical academic reference is Davis~\cite{davis2002}, who defines CBP
as ``planning, under uncertainty, to provide capabilities suitable for a
wide range of modern-day challenges.'' The operational implementations
are the NATO Bi-SC Capability Codes and Capability
Statements~\cite{nato_bisc2016}, which articulate capability requirements
along the DOTMLPFI axes (Doctrine, Organisation, Training, Mat\'{e}riel,
Leadership, Personnel, Facilities, Interoperability), and the US
Department of Defense Joint Capabilities Integration and Development
System (JCIDS)~\cite{jcids2021}, which chains Capabilities-Based
Assessment through Initial Capabilities Document, Capability Development
Document, and Capability Production Document with DOTMLPF-P change
recommendations.

The terminological distinctions within CBP are precise and
methodologically significant. A \textit{capability} is an
outcome-oriented, solution-agnostic ability to achieve a desired effect.
A \textit{technology} is one of the means (the ``M'' in DOTMLPF). A
\textit{priority} is a ranking applied to capabilities or gaps. A
\textit{focus area} is a policy-level thematic cluster at a higher
abstraction level than a capability. A \textit{requirement} is a testable
performance attribute that implements a capability. Davis's central
methodological warning---avoid ``jumping to the solution''---is the logic
that justifies separating capability identification from capability
analysis in any structured planning pipeline.

CBP has crossed into cybersecurity only partially. Lehto and
Limn\'{e}ll~\cite{lehto2022} apply DOTMLPF-I in cyber defence contexts.
AlDaajeh and Alrabaee~\cite{aldaajeh2024} synthesise seven national
cybersecurity strategy axes (legal, technical, organisational,
capacity-building, cooperation, economic, socio-cultural) and
AlDaajeh~\textit{et al.}~\cite{aldaajeh2022} map national cyber
strategies onto capability objectives. Curtis and
Mehravari~\cite{curtis2015} bridge CERT resilience management with
capability maturity. However, no widely cited publication imports the
full Davis/JCIDS analytical machinery into civilian cybersecurity
roadmapping. In established cyber frameworks, ``capability'' is used
loosely: the NIST CSF~2.0~\cite{nist_csf2} employs \textit{Functions,
Categories, Subcategories, Outcomes}; the NICE
framework~\cite{nist_nice} uses \textit{Tasks, Knowledge, Skills, Work
Roles}; the MITRE ATT\&CK framework is an adversary-behaviour taxonomy
(\textit{Tactics, Techniques, Procedures}), not a defender-capability
taxonomy; and ENISA's National Capabilities Assessment
Framework~\cite{enisa_ncaf} uses ``capability'' closer to a maturity
domain than to Davis's outcome-oriented definition. The present work
adopts Davis's definition explicitly, grounding the terminology in the
CBP tradition to achieve greater precision than prior EU cyber roadmaps.

% -----------------------------------------------------------------
% 3.3 MULTI-DIMENSIONAL ASSESSMENT
% -----------------------------------------------------------------
\subsection{Multi-Dimensional Capability Assessment Frameworks}
\label{sec:rw_multidim}

No single peer-reviewed framework combines technical readiness,
regulatory alignment, market dynamics, and sovereignty assessment in a
unified per-capability template. The relevant contributions are
distributed across four disciplines.

The defence tradition supplies DOTMLPF-P and DOTMLPFI, providing
eight-axis internal decomposition of capability solutions but lacking
external market and regulatory axes. The strategic management tradition
supplies PESTEL (Political, Economic, Social, Technological,
Environmental, Legal), recently applied to cybersecurity by Morgado
\textit{et al.}~\cite{morgado2021}. The engineering tradition supplies
the readiness-level suite: TRL~\cite{mankins1995}, MRL, IRL, and
SRL~\cite{sauser2008}, reviewed systematically by Olechowski, Eppinger,
and Joglekar~\cite{olechowski2015}. These readiness scales are strictly
technical; they do not natively encode regulatory, market, or sovereignty
dimensions. The political science tradition supplies the digital
sovereignty axis, with canonical references including Pohle and
Thiel~\cite{pohle2020}, Roberts~\textit{et al.}~\cite{roberts2021}
(defining five sovereignty sub-dimensions: data governance, platform
power, digital infrastructure, emerging technology, and cybersecurity),
Celeste~\cite{celeste2023}, and Falkner~\textit{et al.}~\cite{falkner2024}.
None of these constitutes a per-capability profiling template.

Cybersecurity maturity models come closest to multi-dimensional
profiling. The Oxford Global Cyber Security Capacity Centre Maturity
Model for Nations (GCSCC CMM)~\cite{gcscc_cmm} uses five dimensions,
24~factors, and five maturity stages (start-up through dynamic).
Structurally, it follows a four-step logical pattern---framework
identification, capability enumeration, profiling,
phased progression---that closely parallels the pipeline proposed in this
paper, but it is applied to national-level capacity assessment rather
than technical cybersecurity capabilities and does not anchor to
architectural reference models. The DOE C2M2~\cite{doe_c2m2} employs
ten domains with approximately 356 practices scored at MIL~0--3. The
NIST CSF~2.0~\cite{nist_csf2} uses six Functions, Categories,
Subcategories, and four implementation Tiers. The ITU Global
Cybersecurity Index uses five pillars (legal, technical, organisational,
capacity development, cooperation). AlDaajeh and
Alrabaee~\cite{aldaajeh2024} synthesise seven national strategy axes
but lack sovereignty and readiness dimensions.

The gap is clear: GCSCC covers technical, regulatory, societal,
educational, and strategic aspects, but sovereignty and market dynamics
are only implicit. ENISA's National Capabilities Assessment
Framework~\cite{enisa_ncaf} covers regulatory and sectoral maturity but
omits TRL and market axes. Roberts~\textit{et al.}~\cite{roberts2021}
add sovereignty but at the policy-discourse level, not as a
per-capability metric. A profiling template that fuses technical
readiness, regulatory alignment, market feasibility, and an explicit
digital sovereignty axis within a single per-capability instrument---as
proposed in the 7-Dimensional Capability Profile of
Section~\ref{sec:stage3}---has no direct precedent in the surveyed
literature.

% -----------------------------------------------------------------
% 3.4 EUROPEAN CYBERSECURITY ROADMAP PRECEDENTS
% -----------------------------------------------------------------
\subsection{European Cybersecurity Roadmap Precedents}
\label{sec:rw_eu}

Every major European cybersecurity roadmap deliverable produced under the
Horizon~2020 Cybersecurity Competence Network pilots has been reviewed.
A consistent observation emerges: none cites the canonical TRM
literature (Phaal, Garcia \& Bray, Kostoff) directly, and none employs
quantitative per-capability maturity metrics.

The sole EU precedent that formally grounds cybersecurity roadmapping in
TRM theory is the CyberROAD FP7 project. Ariu~\textit{et al.}~\cite{ariu2016}
cite Phaal, Probert, Carvalho, and the EIRMA formulation, proposing a
co-evolutionary model of technology, threat, and policy trajectories.
CyberROAD omits architecture-anchored capability derivation,
regulatory and sovereignty profiling axes, and quantitative maturity
progression. It is the single most important methodological predecessor
to the present work.

The four H2020 Competence Network pilots each contribute distinct
methodological elements. CONCORDIA~\cite{concordia} (Deliverable~D4.4,
2022) organises its roadmap across six dimensions (Research \&
Innovation; Education \& Skills; Legal \& Policy; Economics \&
Investments; Certification \& Standardisation; Community Building) with
a per-dimension narrative pipeline from threat landscape through
challenge identification to phased recommendations across three time
horizons. It offers the strongest sovereignty-aware multi-dimensional
framing among the pilots but employs no quantitative maturity metrics.
SPARTA~\cite{sparta_d34} (Deliverables~D3.1--D3.5, 2019--2022) structures
challenges against the JRC Cybersecurity Taxonomy (domains $\times$
sectors) with a per-challenge template (Final Goal, EU Benefits, JRC
Domain, JRC Sector, Emerging Technology relation). The taxonomy-grounded
derivation is a strength, but quantitative maturity tracking is absent.
CyberSec4Europe~\cite{cybersec4europe_d45} (Deliverable~D4.5, 2022) is
the strongest multi-dimensional precedent: a structured per-vertical
schema across seven application verticals profiles each challenge across
SWOT analysis combined with Digital Sovereignty, Green Deal, and Digital
Decade Strategy alignment, phased into three time horizons. Its
methodology section explicitly paraphrases Phaal's formulation of
roadmaps as ``structured visual representations.''
ECHO~\cite{echo_d43} (Deliverables~D4.3 and D4.10) takes a
methodological outlier path by mapping SCRUM onto roadmap
development---EPICs, User Stories, acceptance criteria---and adds
Analytic Hierarchy Process (AHP)-based multi-sector assessment.
Quantitative prioritisation is present, but no unified phased-maturity
matrix is produced.

Beyond the Competence Network, the European Cyber Security Organisation
(ECSO) Strategic Research and Innovation Agenda~\cite{ecso_sria} uses
SWOT analysis across thematic pillars without per-priority
multi-dimensional templates or maturity metrics. ENISA's Threat
Landscape Methodology~\cite{enisa_ctl} supplies the most explicit
multi-stage analytical pipeline in EU cybersecurity---six steps
(Direction $\rightarrow$ Collection $\rightarrow$ Processing
$\rightarrow$ Analysis \& Production $\rightarrow$ Dissemination
$\rightarrow$ Feedback)---but adapts the classical intelligence cycle,
not TRM, and is threat-centric rather than capability-centric. The ECCC
Strategic Agenda~\cite{eccc} is a policy strategy rather than an
operational roadmap methodology.

Table~\ref{tab:eu_terminology} summarises the terminological
heterogeneity across EU cybersecurity roadmaps, illustrating that
``capability'' is almost never employed as the formal primary unit of
analysis.

\begin{table}[!ht]
\centering
\caption{Primary analytical units in EU cybersecurity roadmaps}
\label{tab:eu_terminology}
\footnotesize
\begin{tabular}{p{2.8cm} p{2.5cm} p{2.2cm}}
\toprule
\textbf{Project} & \textbf{Primary Unit} & \textbf{Grouping} \\
\midrule
CONCORDIA & Recommendations, challenges & Dimensions \\
ECSO & Priority areas, challenges & Pillars, WGs \\
ENISA & Threats, incidents & Landscapes \\
SPARTA & Challenges & JRC domains \\
CyberSec4Europe & Research challenges & Verticals \\
ECHO & User Stories & EPICs \\
ECCC & Investment priorities & Strategic pillars \\
\bottomrule
\end{tabular}
\end{table}

% -----------------------------------------------------------------
% 3.5 GAP ANALYSIS AND POSITIONING
% -----------------------------------------------------------------
\subsection{Gap Analysis and Methodological Positioning}
\label{sec:rw_gap}

The four research streams surveyed above are consolidated against the
four pipeline elements that define the present methodology:
(a)~architecture-anchored capability derivation, (b)~capability
enumeration from architecture, (c)~multi-dimensional profiling covering
technical, regulatory, market, and sovereignty axes, and (d)~phased
progression matrix with quantitative maturity metrics.
Table~\ref{tab:gap_analysis} presents the assessment.

\begin{table*}[!ht]
\centering
\caption{Gap analysis: existing frameworks assessed against the four
pipeline elements of the proposed methodology. Entries indicate coverage:
Yes, Partial, Weak, or No.}
\label{tab:gap_analysis}
\footnotesize
\begin{tabular}{p{3.8cm} c c c c c}
\toprule
\textbf{Framework} &
\textbf{(a) Arch.} &
\textbf{(b) Enum.} &
\textbf{(c) Multi-dim.} &
\textbf{(d) Quant.} &
\textbf{All 4?} \\
\midrule
CyberROAD~\cite{ariu2016}
    & Partial & Yes & Partial & No & \textbf{No} \\
NIST CSF 2.0~\cite{nist_csf2}
    & Yes & Partial & Weak & Partial & \textbf{No} \\
ENISA NCAF~\cite{enisa_ncaf}
    & Partial & Yes & Partial & Yes & \textbf{No} \\
DOE C2M2~\cite{doe_c2m2}
    & No & Yes & No & Yes & \textbf{No} \\
Phaal T-Plan~\cite{phaal2004}
    & Generic & Yes & Partial & Possible & \textbf{No} \\
GCSCC CMM~\cite{gcscc_cmm}
    & No & Yes & Partial & Yes & \textbf{No} \\
CONCORDIA D4.4~\cite{concordia}
    & Yes & Yes & Partial & No & \textbf{No} \\
SPARTA D3.4~\cite{sparta_d34}
    & Yes & Yes & Partial & No & \textbf{No} \\
CyberSec4Europe D4.5~\cite{cybersec4europe_d45}
    & Yes & Yes & Strong & No & \textbf{No} \\
ECHO D4.3~\cite{echo_d43}
    & Yes & Yes & Partial & Partial & \textbf{No} \\
\textbf{This work}
    & \textbf{Yes} & \textbf{Yes} & \textbf{Yes} & \textbf{Yes}
    & \textbf{Yes} \\
\bottomrule
\end{tabular}
\end{table*}

No existing framework integrates all four pipeline elements.
CyberSec4Europe D4.5 is the closest precedent on multi-dimensional
profiling; ECHO is the closest on quantitative prioritisation via AHP;
ENISA's Threat Landscape methodology is the closest on explicit
multi-stage analytical pipelines; CONCORDIA is the closest on
sovereignty-driven framing; and the GCSCC CMM is the closest structural
analogue but operates as a national-capacity maturity model rather than
a technical roadmap anchored to architectural reference models.

The present methodology occupies this gap by fusing three traditions:
the Cambridge/Sandia TRM tradition~\cite{phaal2004,garcia1997} for
phased planning structure and stage-gate logic, the Davis/JCIDS
capability-based planning tradition~\cite{davis2002} for formal
capability definition and architecture-to-capability derivation, and
cybersecurity maturity models~\cite{nist_csf2,gcscc_cmm,doe_c2m2}
for quantitative progression metrics. The novelty is defensible on
three independent grounds: (i)~no canonical TRM framework separates
capability identification from capability analysis as named, gated
stages; (ii)~no EU cybersecurity roadmap employs quantitative
per-capability maturity metrics anchored to an architectural reference
model; and (iii)~no peer-reviewed framework profiles capabilities
simultaneously across technical readiness, regulatory alignment, market
dynamics, and digital sovereignty in a single template.


